\section{Introduction}
\label{sec:intro}

Automatic solar-filament segmentation requires detecting thin, low-contrast structures
in full-disc images. The task is complicated by the strong imbalance between foreground
and background, the variation in filament shape and size, and the availability of
multiple annotations for the same observation. A useful prediction must therefore
identify filament pixels without producing too many spurious or fragmented regions.

This work develops a reproducible pipeline for image segmentation and subsequent
instance extraction. Its main model is a U-Net with a pre-trained ResNet18 encoder,
chosen to combine semantic information with spatial detail. Mean Dice is the primary
metric, while Panoptic Quality (PQ) also evaluates instance recognition. GPU-memory use
and execution time are measured because the solution must be trainable and usable with
limited resources.

During development, \textbf{22 experimental configurations were run}. The report focuses
on the eight numbered experiments that determined the main design choices. It also compares the
full-resolution U1 and U2 variants and the three YOLO detector experiments. This
organisation reconstructs the path from the baseline to the delivered model while also
showing alternatives that improved PQ at a higher computational cost.

The following sections define the metrics and validation protocol, describe the
architecture and training procedure, and discuss the most significant experiments. The
final part analyses prediction errors, computational cost and the main limitations of
the solution.
